% Generated draft: requires mathematical and source review.
\documentclass[11pt]{article}
\usepackage[margin=25mm]{geometry}
\usepackage{fontspec}
\setmainfont{Latin Modern Roman}
\usepackage{amsmath,amssymb,mathtools,unicode-math}
\setmathfont{Latin Modern Math}
\usepackage{enumitem,needspace,xurl,hyperref}
\setlength{\parindent}{0pt}
\setlength{\parskip}{5pt}
\setlength{\emergencystretch}{3em}
\newcommand{\N}{\mathbb{N}}
\newcommand{\Z}{\mathbb{Z}}
\newcommand{\Q}{\mathbb{Q}}
\newcommand{\R}{\mathbb{R}}
\newcommand{\C}{\mathbb{C}}
\newcommand{\F}{\mathbb{F}}
\newcommand{\A}{\mathbb{A}}
\newcommand{\PP}{\mathbb{P}}
\newcommand{\E}{\mathbb{E}}
\newcommand{\eps}{\varepsilon}
\newcommand{\dd}{\,\mathrm{d}}
\begin{document}
\section*{UnsolvedMath ID: 20000753}
\textbf{Source problem code:} catalogue entry
\subsection{Definitions and mathematical statement}
J.B. Conrey asks whether a suitable Elliott--Halberstam (EH) statement for prime $k$-tuples can imply prime $(k+1)$-tuples. T.D. Wooley and others emphasize that the correct tuple-EH hypothesis is itself unspecified. With $\mathbf1_{\mathcal P}$ the prime indicator, define
\[
 E(x,q;a,b)=\sum_{1\leq n\leq x/q}\mathbf1_{\mathcal P}(qn+a)\mathbf1_{\mathcal P}(qn+b)-\operatorname{MT}(x,q;a,b).
\]
The notes mention Balog--Kawada--Mikawa results of the schematic form
\[
 \sum_{q\leq Q}\max_{(a,q)=1}\sum_{\substack{1\leq b\leq q\\(b,q)=1}}
       |E(x,q;a,b)|\ll_A x^2(\log x)^{-A}.
\]
R.C. Vaughan proposes the different possible hypothesis
\[
 \sum_{q\leq Q}\max_{(a,q)=1}
 \left|\sum_{\substack{1\leq b\leq B\\(b,q)=1}}E(x,q;a,b)\right|
 \ll_A Bx(\log x)^{-A}.
\]
The position of the absolute value is different. The extracted comparison signs are damaged; the displays restore the intended bound notation only schematically. No main term, allowed $Q,B$ range, representative convention for $a$, or uniform quantifiers are supplied. J. Pintz separately asks whether EH for twins can yield two boundedly close twin-prime pairs.

The aid's finite transfer criterion makes the extra correlations explicit. A finite set $\mathcal H$ of $k\geq1$ distinct integer shifts is admissible if it occupies fewer than $p$ residue classes modulo every prime $p$. Choose distinct extension shifts in a finite set $\mathcal T$ disjoint from $\mathcal H$, with every one- and two-shift enlargement admissible. The aid reports such sets always exist: for instance, with $h_0\in\mathcal H$ and $M=\prod_{p\leq k+2}p$, choose distinct new shifts congruent to $h_0$ modulo $M$. The larger primes cannot be covered by at most $k+2$ shifts.

On a fixed finite interval of integers $\mathcal I$, put
\[
 A(n)=\prod_{h\in\mathcal H}\mathbf1_{\mathcal P}(n+h),\qquad
 R(n)=\sum_{t\in\mathcal T}\mathbf1_{\mathcal P}(n+t),
\]
\[
 S_0=\sum_{n\in\mathcal I}A(n),\quad
 S_1=\sum_{t\in\mathcal T}\sum_{n\in\mathcal I}A(n)\mathbf1_{\mathcal P}(n+t),
\]
\[
 S_2=S_1+\sum_{\substack{t,u\in\mathcal T\\t\neq u}}
 \sum_{n\in\mathcal I}A(n)\mathbf1_{\mathcal P}(n+t)\mathbf1_{\mathcal P}(n+u).
\]
These are exactly $\sum A$, $\sum AR$ and $\sum AR^2$. If $S_1>0$, the reported moment transfer gives
\[
 \#\{n\in\mathcal I:A(n)=1,\ R(n)>0\}\geq S_1^2/S_2.
\]
More generally, for $0\leq\theta<1$, the usual Paley--Zygmund form bounds the count with $R(n)>\theta S_1/S_0$ below by $(1-\theta)^2S_1^2/S_2$. Thus usable quantitative input involves aggregate $(k+1)$- and $(k+2)$-point correlations, not just marginal progression distribution of $A(n)$.

The aid reports an abstract obstruction using disjoint prime-density zero-one sequences, each with marginal progression discrepancies smaller than $N/(\log N)^C$ for every fixed $C$ and every fixed distribution level below one. Their disjointness shows that marginal distribution quality alone does not force overlap or a shifted tuple extension. The summary does not specify the model's exact progression main terms; it is not a counterexample to any precisely formulated assertion about the actual primes, which have additional structure.

For Pintz's four-distinct-shift twin-pair pattern $\{0,2,d,d+2\}$, with positive separation $d>2$, the reported local-admissibility criterion is exactly $d\geq6$ and $6\mid d$. An infinite bounded-separation twin-pair conclusion reduces to a finite disjunction of fixed admissible patterns $\{0,2,6m,6m+2\}$, with the separation bound fixed. Admissibility and this reduction alone do not prove such prime patterns occur infinitely often.
\subsection{Short English statement}
The original question leaves tuple EH undefined. The aid shows that a moment argument for extending tuples needs higher shifted correlations, and gives an abstract obstruction to inferring overlap from marginal progression data alone. It also identifies the admissible fixed patterns for two nearby twin-prime pairs, without proving their occurrence.
\subsection{Sources}
\begin{itemize}
\item Local UnsolvedMath registry, permanent ID $20000753$, source code \texttt{catalogue entry}, supplied full source record and research summary.
\item Angel Kumchev's notes, December 2005 ARCC workshop \emph{Gaps Between Primes}, page 5, Problem 18, ``From $k$-tuples to $(k+1)$-tuples'', as identified in the supplied verification: \url{https://aimath.org/WWN/primegaps/WorkshopProblems.pdf}. Canonical local record \texttt{aim-analytic-number-theory-notes.json}, \texttt{notes[116]}. No PDF was retrieved.
\item Original attributions: J.B. Conrey; T.D. Wooley and others; Balog, Kawada and Mikawa; R.C. Vaughan; J. Pintz. No full bibliographic entries for the mentioned results or a precise tuple-EH conjecture are supplied.
\item Research attempt 1: partial result; problem partially solved at the recorded run; supplied rigor label ``proved'', novelty confidence ``low''. The machine-generated aid is not peer reviewed and requires independent verification before citation.
\end{itemize}
\end{document}
